\subsection{Projective modules of rank 1}

THEOREM 3. Let A be a ring and M a finitely generated A-module.

\begin{enumerate}
\def\labelenumi{(\roman{enumi})}
\item
  If there exists an A-module N such that \(\mathbf{M} \otimes_{\mathbf{A}} \mathbf{N}\) is isomorphic to A, the module M is projective of rank 1.
\item
  Conversely, if \(\mathbf{M}\) is projective of rank 1 and \(\mathbf{M}^*\) is the dual of \(\mathbf{M}\) , the canonical homomorphism \(u: \mathbf{M} \otimes_{\mathbf{A}} \mathbf{M}^* \to \mathbf{A}\) corresponding to the canonical bilinear form \((x, x^*) \to \langle x, x^* \rangle\) on \(\mathbf{M} \times \mathbf{M}^*\) (Algebra, Chapter II, §2, no. 3) is bijective.
\item
  It is required to prove that, for every maximal ideal m of A, the \(A_{m}\) -module \(M_{m}\) is free of rank 1 (Theorem 2 (b)); we are free to replace A by \(A_{m}\) and hence may assume that A is a local ring ( \(\S\) 2, no. 7, Proposition 18). Let k = A/m. The isomorphism v: \(M \otimes_{A} N \to A\) defines an isomorphism
\end{enumerate}

\[
v \otimes 1 _ {k} \colon (\mathbf {M} / \mathfrak {m M}) \otimes_ {k} (\mathbf {N} / \mathfrak {m N}) \rightarrow k
\]

as the rank over \(k\) of (M/mM) \(\otimes_k\) (N/mN) is the product of the ranks of M/mM and N/mN, these latter are necessarily equal to 1, in other words M/mM is monogenous. It follows that M is monogenous (§ 3, no. 2, Corollary 2 to Proposition 4); on the other hand, the annihilator of M also annihilates M \(\otimes_A\) N and hence is zero, which proves that M is isomorphic to A.

\begin{enumerate}
\def\labelenumi{(\roman{enumi})}
\setcounter{enumi}{1}
\tightlist
\item
  It is sufficient to prove that, for every maximal ideal m of A, \(u_{m}\) is an isomorphism ( \(\S 3\) , no. 3, Theorem 1). As M is finitely presented (Chapter I, \(\S 2\) , no. 8, Lemma 8), \((\mathbf{M}^{*})_{\mathfrak{m}}\) is canonically identified with the dual \((\mathbf{M}_{\mathfrak{m}})^{*}\) ( \(\S 2\) , no. 7, Proposition 19) and, as \(M_{m}\) is free of rank 1 like its dual \((\mathbf{M}_{\mathfrak{m}})^{*}\) , clearly the canonical homomorphism \(u_{\mathfrak{m}}: (\mathbf{M}_{\mathfrak{m}}) \otimes_{\mathbf{A}\mathfrak{m}} (\mathbf{M}_{\mathfrak{m}})^{*} \to \mathbf{A}_{\mathfrak{m}}\) is bijective, which completes the proof.
\end{enumerate}

Remark (1). If \(M\) is projective of rank 1 and \(\mathbf{N}\) is such that \(\mathbf{M} \otimes_{\mathbf{A}} \mathbf{N}\) is isomorphic to \(\mathbf{A}\) , then \(\mathbf{N}\) is isomorphic to \(\mathbf{M}^*\) : there are isomorphisms

\[
\mathbf {N} \rightarrow \mathbf {N} \otimes \mathbf {A} \rightarrow \mathbf {N} \otimes \mathbf {M} \otimes \mathbf {M} ^ {*} \rightarrow \mathbf {A} \otimes \mathbf {M} ^ {*} \rightarrow \mathbf {M} ^ {*}.
\]

PROPOSITION 6. Let M and N be projective A-modules of rank 1. Then M ⊗\_\{A\} N, Hom\_\{A\}(M, N) and the dual M* of M are projective of rank 1.

This follows immediately from formulae (2), (3) and (4).

Let us now note that every finitely generated A-module is isomorphic to a quotient module of \(L = A^{(N)}\) ; we may therefore speak of the set \(F(A)\) of classes of finitely generated A-modules with respect to isomorphism (Set Theory, Chapter I, §6, no. 9); we denote by \(P(A)\) the subset of \(F(A)\) consisting of the classes of projective A-modules of rank 1 and by \(\operatorname{cl}(M)\) the image in \(P(A)\) of a projective A-module M of rank 1. It is immediate that, for two projective A-modules M,

N of rank 1, \(\operatorname{cl}(\mathbf{M} \otimes_{\mathbf{A}} \mathbf{N})\) depends only on \(\operatorname{cl}(\mathbf{M})\) and \(\operatorname{cl}(\mathbf{N})\) ; as definition we set

\[
\operatorname{cl} (\mathbf {M}) + \operatorname{cl} (\mathbf {N}) = \operatorname{cl} (\mathbf {M} \otimes_ {\mathbf {A}} \mathbf {N})\tag{6}
\]

and an internal law of composition is thus defined on \(\pmb{P}(\mathbf{A})\) .

PROPOSITION 7. The set P(A) of classes of projective A-modules of rank 1, with the law of composition (6), is a commutative group. If M is a projective A-module of rank 1 and M* is its dual, then

\[
\operatorname{cl} (\mathbf {M} ^ {*}) = - \operatorname{cl} (\mathbf {M}) \quad a n d \quad \operatorname{cl} (\mathbf {A}) = 0.\tag{7}
\]

The associativity and commutativity of the tensor product show that the law of composition (6) is associative and commutative; the isomorphism between \(A \otimes_{A} M\) and M prove that \(\operatorname{cl}(A)\) is the identity element under this law and, by virtue of Theorem 3, \(\operatorname{cl}(M) + \operatorname{cl}(M^{*}) = \operatorname{cl}(A)\) , whence the proposition.

Let B be a commutative A-algebra and M a projective A-module of rank 1; then \(\mathbf{M}_{(\mathbf{B})} = \mathbf{B} \otimes_{\mathbf{A}} \mathbf{M}\) is a projective B-module of rank 1 (no. 3, Proposition 4). Then there exists a mapping called canonical \(\phi: P(\mathbf{A}) \to P(\mathbf{B})\) such that

\[
\phi (\operatorname{cl} (M)) = \operatorname{cl} (M _ {(B)}).\tag{8}
\]

The equation \(\mathbf{M}_{(\mathbf{B})}\otimes_{\mathbf{B}}\mathbf{N}_{(\mathbf{B})}=(\mathbf{M}\otimes_{\mathbf{A}}\mathbf{N})_{(\mathbf{B})}\) for two A-modules M, N proves that the mapping \(\phi\) is a commutative group homomorphism.

*Remark (2). Condition (e) of Theorem 1 (equivalent to the fact that P is projective and finitely generated) may also be expressed by saying that the sheaf of modules \(\tilde{\mathbf{P}}\) over \(\mathbf{X} = \operatorname{Spec}(\mathbf{A})\) associated (*) with P is locally free and of finite type and may therefore be interpreted as the sheaf of sections of a vector bundle over X. Conversely, every vector bundle over X arises from a finitely generated projective module, which is determined to within a unique isomorphism; the projective modules of rank \(n\) thus correspond to the vector bundles all of whose fibres have dimension \(n\) . In particular, the vector bundles of rank 1 correspond to the projective modules of rank 1. If we denote by \(\mathcal{O}_{\mathbf{x}}\) the structure sheaf \(\tilde{\mathbf{A}}\) and by \(\mathcal{O}_{\mathbf{x}}^{*}\) the sheaf of units of \(\mathcal{O}_{\mathbf{x}}\) (whose sections over an open set U of X are the invertible elements of the ring of sections of \(\mathcal{O}_{\mathbf{x}}\) over U), it follows that the group \(P(\mathbf{A})\) is isomorphic to the first cohomology group \(\mathrm{H}^1 (\mathrm{X},\mathcal{O}_\mathrm{x}^*)\) .

